% Proved after full independent root and companion audits; exact candidate preserved.
\subsection{Iterated actual high quotients and a third-index interval}

For the local assertions let \(K\) be complete discretely valued
of mixed characteristic \((0,p)\), with algebraically closed
residue field, \(p\ge5\), \(\mu_p\subset K\), and absolute
ramification \(e\) prime to \(p\). Set
\[
 q=p^2,\qquad N=q-1,\qquad s=e(p+1).
 \tag{B-RD74.1}
\]
Use the actual common good tame field \(L/K\) of degree \(N\),
so \(v_L(p)=eN\). The two-slope parameter and high-line
character are those of Lemma~\ref{lem:B-RD69-rules}:
\[
 0<B<pe,\qquad \alpha=s-B,\qquad\beta=B/p,
 \qquad \lambda_H=\zeta^{c-B}.
 \tag{B-RD74.2}
\]
There is no square-class restriction on the finite determinant.

\begin{lemma}[Actual levels in a sequence of high quotients]
\label{lem:B-RD74-actual-levels}
\status{Proved; complete independent root and companion proof and exact-code audits passed.}
Let \(n\ge2\) and \(0\le j\le n-1\). Suppose actual
\(p^n\)-torsion modules of
\(M_1,M_2\) are identified, and at each of the first \(j\)
stages the chosen high lines are identified. Write
\(M_i^{(0)}=M_i\) and
\(M_i^{(j+1)}=M_i^{(j)}/H_i^{(j)}\).
After \(j\) quotients the induced comparison is actual
\(p^{n-j}\)-torsion. Each quotient is the actual good formal
quotient over \(\mathcal O_L\), with its inherited generic
semilinear descent and the same finite special character.

If a stage has divisible parameter \(D\), then \(D\ne e\).
For \(D<e\) the next group has parameter \(pD\) and its
first-level module has two distinct stable lines, hence is
\(L\)-rational. For \(D>e\) the next parameter \(s-D\) is
prime to \(p\), and its first-level division field over \(L\)
is wild of degree \(p\).
\end{lemma}
\begin{proof}
If the current level is \(p^r\), with \(r\ge2\), the exact
generic finite-group construction is
\[
 (M/H)[p^{r-1}]
  =\{X\in M[p^r]:p^{r-1}X\in H\}/H.
 \tag{B-RD74.3}
\]
Indeed a lift of a point killed by \(p^{r-1}\) has its
\(p^{r-1}\)-multiple in \(H\), hence is killed by \(p^r\).
Conversely this condition is sufficient, and the kernel is
exactly \(H\). Applying the actual torsion isomorphism to this
subgroup and quotient gives the asserted comparison. Iteration
loses one level at every stage, including at the final stage,
which retains only \(p\)-torsion.

The finite-flat closure of each stable generic line is stable
under the actual semilinear descent action. The good formal
quotient, quotient coordinate, and effective generic descent
are exactly those proved in
Lemma~\ref{lem:B-RD69-quotients}; its special isogeny preserves
the rational finite special character. The total kernel after
\(j\) stages has rank \(p^j\), is killed by \(p^j\), and its
preimage construction is contained in the original
\(M[p^n]\). Thus no comparison of higher truncation than
\(p^{n-j}\) is being asserted.

When \(p\mid D\), equality \(D=e\) is excluded by \(p\nmid e\).
The exact actual high-quotient parameter is \(pD\) for \(D<e\)
and \(s-D\) for \(D>e\). In the first case the intrinsic
image of \(M[p]\) in \((M/H)[p]\) is a stable low line,
distinct from its stable high line, by
Lemma~\ref{lem:B-RD71-marked-low}. Both characters are
\(\mathbf F_p^\times\)-valued, so the degree-\(N\) tame field
kills them. This proves actual first-level rationality.
In the second case \(p\nmid(s-D)\). The low valuation
\((s-D)/p\) is nonintegral over \(L\), whereas the high
character and determinant are trivial there. The residual
action is therefore the nontrivial upper-unipotent group of
order \(p\), as in Lemma~\ref{lem:B-RD71-first-level}.
This proves the stated tame/wild alternatives.
\end{proof}

\begin{proposition}[Uniform iteration bound for whole formal groups]
\label{prop:B-RD74-formal}
\status{Proved; complete independent root and companion proof and exact-code audits passed.}
Let \(n\ge3\) and \(e<p^n\). An isomorphism of actual
\(p^n\)-torsion determines the finite special inertia character
among whole potentially supersingular height-two groups.
\end{proposition}
\begin{proof}
The actual \(p^3\)-comparison already available handles both
one-slope polygons and excludes a mixed one-slope/two-slope
comparison by Proposition~\ref{prop:B-RD69-one-slope}, whose
scope has no upper bound on \(e\).
For two two-slope groups, nontrivial residual wild inertia
identifies the high line as its unique fixed line. If the
residual action is tame, Lemma~\ref{lem:B-RD73-high-branch}
gives actual \(L\)-rational first-level modules, and
Proposition~\ref{prop:B-RD73-high-recovery} identifies the
high lines from the available actual \(p^2\)-comparison.
Thus in every two-slope case the original high lines agree.

Write the original parameters \(B_i\) and tangent exponents
\(c_i\). Their common finite determinant and high character
give
\[
 c_1\equiv c_2\pmod{p-1},\qquad
 c_1-B_1\equiv c_2-B_2\pmod N.
 \tag{B-RD74.4}
\]
If at least one \(B_i\) is prime to \(p\), the general
comparison rules of Lemma~\ref{lem:B-RD69-rules}, using only
the available \(p^2\)-torsion, recover the character or
exclude opposite divisibility. In particular this settles
the possible original equality \(B_i=e\).
Assume henceforth that both \(B_i\) are divisible by \(p\).
If both exceed \(e\), their first high quotients have prime
parameters \(s-B_i\). Their actual first-level comparison
recovers the character by Lemma~\ref{lem:B-RD71-first-level},
and the same holds for the original groups by isogeny
invariance. If they straddle \(e\), the next first-level
modules are tame and wild respectively by the preceding
lemma, which is impossible. Thus the only remaining case is
\[
 B_i=pb_i<e,\qquad b_i>0.
 \tag{B-RD74.5}
\]

We now give the complete induction, including its level count.
As long as the previous parameters were below \(e\), the
actual group after \(j\) high quotients has parameter
\[
 D_i^{(j)}=p^j B_i=p^{j+1}b_i,
 \qquad 0\le j\le n-2,
 \tag{B-RD74.6}
\]
and its available actual comparison has level \(p^{n-j}\).
For \(j=0\) the high lines were already identified. For
\(1\le j\le n-2\), the preceding small quotient supplies
the intrinsic stable low line in the current group. Its
first-level module is \(L\)-rational. Since \(n-j\ge2\),
Proposition~\ref{prop:B-RD73-high-recovery} identifies the
current high lines, even for scalar residual modules.
We may therefore take the next actual quotient, leaving
level \(p^{n-j-1}\), which is at least \(p\).

Each \(D_i^{(j)}\) is divisible by \(p\), so it cannot equal
\(e\). If both are greater than \(e\), the next parameters
\(s-D_i^{(j)}\) are prime to \(p\); their actual first-level
comparison recovers equality and the finite characters, by
Lemma~\ref{lem:B-RD71-first-level}. If they lie on opposite
sides of \(e\), the next first-level representations have
trivial and nontrivial wild inertia, by
Lemma~\ref{lem:B-RD74-actual-levels}. This is impossible.
These two conclusions require only the \(p\)-torsion that is
still present even at \(j=n-2\).

The only way to continue through every stage is that both
parameters remain below \(e\), including at \(j=n-2\).
In that case
\[
 0<b_i<\frac{e}{p^{n-1}}<p,
 \qquad b_i\in\{1,\ldots,p-1\}.
 \tag{B-RD74.7}
\]
The common characters in (B-RD74.4) imply
\(B_1\equiv B_2\pmod{p-1}\), hence
\(b_1\equiv b_2\pmod{p-1}\). Reduction modulo \(p-1\) is
injective on precisely \(\{1,\ldots,p-1\}\); thus
\(b_1=b_2\). The common original high character then gives
\(c_1=c_2\pmod N\), proving equality of the special
characters. Every branch has been exhausted, and at no point
was a second-level comparison assumed after only the first
level remained.
\end{proof}

\begin{theorem}[A logarithmic upper bound in the tame absolute-index case]
\label{thm:B-RD74-upper}
\status{Proved; complete independent root and companion proof and exact-code audits passed.}
In the full B.4 monodromy setting of height two, with perfect
residue field, \(p\ge5\), \(\mu_p\subset K\), and \(p\nmid e\),
every integer \(n\ge3\) such that \(e<p^n\) satisfies
\[
 \nu_{\min}(p,e,2)\le n-1.
 \tag{B-RD74.8}
\]
In particular, if \(e>p\) and \(k\) is the unique integer
with \(p^k<e<p^{k+1}\), then
\[
 \left\lfloor\frac{k+3}{2}\right\rfloor
 \le \nu_{\min}(p,e,2)\le\max\{2,k\}.
 \tag{B-RD74.9}
\]
\end{theorem}
\begin{proof}
Pass to the completed maximal unramified base as in the
established character-descent reduction. The cyclotomic
character modulo \(p^n\) has order \(p^{n-1}\), because
\(p\nmid e\); the determinant comparison distinguishes total
dimensions \(0,1,2\) and their finite determinant characters.
The height-one filtered cases are determined residually.
For a mixed whole-supersingular/filtered comparison,
Proposition~\ref{prop:B-RD63-mixed} requires
\(e>p^{2n-3}\). But \(n\ge3\) gives \(n\le2n-3\), so
\(e<p^n\le p^{2n-3}\), excluding this possibility.
Two whole formal groups are covered by the preceding
proposition. This proves (B-RD74.8) in the full setting,
rather than only for potentially good whole groups.

For the upper bound in (B-RD74.9), take
\(n=\max\{3,k+1\}\). For the lower bound set
\(r=\lfloor(k+3)/2\rfloor\). Then \(r\ge2\) and
\(2r-3\le k\), so \(e>p^k\ge p^{2r-3}\).
The actual \(p^r\)-torsion lower construction in
Theorem~\ref{thm:B-RD63-arc} therefore gives
\(\nu_{\min}\ge r\). The conversion uses the manuscript
convention that actual \(p^r\)-torsion is its index \(r-1\).
No equality of these two bounds is asserted when they differ.
\end{proof}

\begin{corollary}[A complete third-index interval]
\label{cor:B-RD74-third-index}
\status{Proved; complete independent root and companion proof and exact-code audits passed.}
Under the hypotheses of the preceding theorem,
\[
 p^3<e<p^4\quad\Longrightarrow\quad
 \boxed{\nu_{\min}(p,e,2)=3}.
 \tag{B-RD74.10}
\]
For \(p=5\), this includes every multiple of four from
\(128\) through \(624\) which is prime to five.
\end{corollary}
\begin{proof}
Take \(n=4\) in the upper bound. Since \(e>p^3\),
Theorem~\ref{thm:B-RD63-arc} with \(r=3\) supplies actual
\(p^3\)-torsion counterexamples, so the least index is at
least three. These bounds coincide. The numerical endpoints
are the first and last multiples of four in \((125,625)\).
\end{proof}

\begin{remark}[Scope and remaining interval]
\label{rem:B-RD74-scope}
\status{Proved; complete independent root and companion proof and exact-code audits passed.}
The iteration uses actual generic torsion, actual finite-flat
high quotients, and their generic semilinear descent. The
second line ensuring rationality is the intrinsic image of
the preceding first-level module; no arbitrary complement
or scalar eigenspace is chosen. The exact third-index result
is a sufficient interval, not an ``if and only if'' statement.
For example, \(p^4<e<p^5\) remains bounded only by
\(3\le\nu_{\min}\le4\) here. The arguments do not settle
absolute ramification divisible by \(p\), higher heights,
or the exact least index throughout all larger intervals.
The retained finite supplement is
\path{code/reduction_descent/verify_iterated_high_quotient.py}.
Its final author run passed in \(0.863016332994448\) seconds
at 22:44:19 UTC on 19 November 2026. It exhaustively verifies
four actual rank-two module quotients at \(p=5,7\), levels
two and three, and \(81{,}600\) induced-matrix point identities.
It checks all \(196\) permitted sectors for
\((p,n)=(5,3),(5,4),(7,3)\), including every compatible small
parameter pair, and \(406\) explicitly sampled deeper sectors
for \(p=5,7,11,13\) and \(3\le n\le8\).
The complete and sampled parameter-pair counts are respectively
\(88{,}870\) and \(161{,}887\); among the checked stages,
\(16{,}832\) stop when only the final first level remains.
The program also checks the unordered character orbit under
each isogeny and the precise lower/upper-bound arithmetic.
These are finite algebra checks, not a proof of formal-group
existence or a substitute for the induction.

The first executed program and its successful output are
preserved. A subsequent harmless rewrite replaced division by
multiplication in one integer inequality, and the final exact
integer version was rerun successfully. No executed check
failed. The source chain is the fully audited actual quotient,
Newton, lower-break and mixed-threshold proofs B69/B71/B73/B63;
no additional theorem is inferred from the bounded B73
literature search.
The complete independent root and companion proof/code audits
passed. Their author-output-preserving exact reruns took
\(0.8197192030056613\) and \(0.8659103850077372\) seconds,
respectively. All five assertions, exact candidate bytes,
executed versions, and audit records are retained.
\end{remark}
